% Archivo: sections/2_theory.tex
\section{Marco Teórico}

La dinámica de los sistemas estudiados se rige por sistemas de ecuaciones diferenciales ordinarias no lineales. Para el péndulo simple forzado y amortiguado, la ecuación de movimiento se deriva de la segunda ley de Newton, considerando una fuerza de arrastre proporcional a la velocidad angular y una fuerza impulsora armónica. Definiendo el ángulo $\theta$ y la velocidad angular $\omega$, el sistema se reduce a dos ecuaciones de primer orden:
\begin{equation}
    \dot{\theta} = \omega
\end{equation}
\begin{equation}
    \dot{\omega} = -\frac{g}{l}\sin(\theta) - q\omega + F_D \sin(\Omega t)
\end{equation}

donde $q$ es el coeficiente de amortiguamiento y $F_D$ la amplitud de la fuerza impulsora con frecuencia $\Omega$.


Por otro lado, el péndulo doble constituye un sistema conservativo con dos grados de libertad ($\theta_1, \theta_2$) cuyas ecuaciones de movimiento se obtienen mediante el formalismo de Euler-Lagrange. Dadas las masas $m_1, m_2$ y las longitudes de las varillas $L_1, L_2$, las aceleraciones angulares acopladas vienen dadas por:
\begin{equation}
    \ddot{\theta}_1 = \frac{-g(2m_1 + m_2)\sin\theta_1 - m_2 g \sin(\theta_1 - 2\theta_2) - 2\sin(\theta_1 - \theta_2)m_2\left(\dot{\theta}_2^2 L_2 + \dot{\theta}_1^2 L_1 \cos(\theta_1 - \theta_2)\right)}{L_1 \left(2m_1 + m_2 - m_2 \cos(2\theta_1 - 2\theta_2)\right)}
\end{equation}

\begin{equation}
    \ddot{\theta}_2 = \frac{2\sin(\theta_1 - \theta_2)\left(\dot{\theta}_1^2 L_1 (m_1 + m_2) + g(m_1 + m_2)\cos\theta_1 + \dot{\theta}_2^2 L_2 m_2 \cos(\theta_1 - \theta_2)\right)}{L_2 \left(2m_1 + m_2 - m_2 \cos(2\theta_1 - 2\theta_2)\right)}
\end{equation}

Este sistema exhibe una extrema dependencia a las condiciones iniciales a pesar de su simplicidad estructural \citep{pend_doble}.
 

El modelo de Lorenz surge como una simplificación extrema de las ecuaciones de Navier-Stokes y de conducción de calor para la convección atmosférica \citep{Lorenz}. El sistema resulta en:
\begin{equation}
    \dot{x} = \sigma (y - x)
\end{equation}
\begin{equation}
    \dot{y} = rx - y - xz
\end{equation}
\begin{equation}
    \dot{z} = xy - bz
\end{equation}

donde $\sigma$, $r$ y $b$ son el número de Prandtl, el número de Rayleigh normalizado y un parámetro geométrico, respectivamente.

Para cuantificar la sensibilidad a las condiciones iniciales, sello distintivo del caos, se define el exponente de Lyapunov, $\lambda$. Si consideramos dos trayectorias inicialmente separadas por una distancia infinitesimal $|\delta \theta(0)|$, su separación evoluciona según:
\begin{equation}
    |\delta \theta(t)| \approx |\delta \theta(0)| e^{\lambda t}
\end{equation}

El signo de $\lambda$ determina cómo evoluciona el sistema en el tiempo:

$\bullet$ $\lambda < 0$ (Estable): Las trayectorias se acercan. El sistema pierde energía y termina en un estado predecible, como un punto fijo o una órbita cerrada (ciclo límite).

$\bullet$ $\lambda = 0$ (Marginal): El sistema ni converge ni se aleja exponencialmente. Es el punto crítico de las transiciones o de sistemas con órbitas estables que no cambian de amplitud.


$\bullet$ $\lambda > 0$ (Caótico): Las trayectorias se separan exponencialmente. Un pequeño cambio inicial altera drásticamente el resultado final. Esto genera caos, donde el sistema es determinista pero imposible de predecir a largo plazo, moviéndose en un atractor extraño.